\section{Related Literature, Contribution, and Research Hypotheses}
\label{sec:related-literature}

This study lies at the intersection of cross-sectional empirical asset pricing,
nonlinear statistical forecasting, time-varying market conditions, and the
economics of implementing return predictions. These literatures overlap, but
they address different questions. Evidence that a model predicts the
\emph{aggregate} equity premium is not, by itself, evidence that it ranks
\emph{individual stocks} accurately. Similarly, statistical accuracy,
cross-sectional rank correlation, gross long--short returns, and investable
profitability are separate objects. The contribution of the present study must
be evaluated against all four distinctions.

\subsection{Machine learning and cross-sectional asset pricing}
\label{subsec:literature-ml}

The empirical asset-pricing literature has long studied how the observable
characteristics of securities relate to their subsequent returns.
\cite{FamaFrench1992} identify important cross-sectional associations involving
size and book-to-market, while \cite{JegadeeshTitman1993} document the
profitability of intermediate-horizon momentum strategies.
The common-risk-factor formulation of \cite{FamaFrench1993} motivates a
distinction between a characteristic's predictive association and its relation
to systematic risk. Subsequent research on gross profitability and investment
illustrates the limitations of a specification containing size, value and
momentum alone \cite{NovyMarx2013,FamaFrench2015}. The three characteristics
used here are therefore interpretable and historically established, but should
not be mistaken for a complete information set.

The introduction of flexible prediction methods has shifted attention from
individual characteristic slopes to high-dimensional and potentially nonlinear
maps from firm information to expected returns. \cite{GuKellyXiu2020} provide
a broad comparison of machine-learning methods for empirical asset pricing and
show the importance of allowing nonlinear interactions among predictors.
\cite{FreybergerNeuhierlWeber2020} study characteristic selection and
nonlinearity in a nonparametric framework, while \cite{KozakNagelSantosh2020}
illustrate the value of economically motivated shrinkage in the cross-section.
\cite{KellyPruittSu2019} model characteristics as conditioning information
for time-varying risk exposures. These studies demonstrate why a
regularized linear predictor and a nonlinear boosted-tree model provide
informative benchmarks, but they also establish that neither regularization
nor nonlinear characteristic modeling is a new methodological invention.
In particular, Gu, Kelly and Xiu explicitly include a three-variable
ordinary-least-squares comparator based on size, book-to-market and
momentum, and compare it with regularized linear and gradient-boosted
models. Their approximately 1957--2016 monthly stock panel uses an
1987--2016 out-of-sample period with annual model refits, as opposed
to the monthly expanding refits of the present project. Thus, using
these three characteristics alone cannot be represented as an original
selection of explanatory variables.

Further evidence on flexible factor representations is provided by
\cite{ChenPelgerZhu2024}. The critical survey of
\cite{Bagnara2024} highlights the range of regularization, dimension
reduction, tree and neural-network approaches already applied to asset
pricing. Thus, comparing Elastic Net with XGBoost should be framed as a
controlled empirical comparison of two established forecasting families,
not as the introduction of a new class of forecasting model.

The present design deliberately keeps the characteristic set small.
The frozen predictors are log market equity, book-to-market, and
momentum 12--2, and the same sample and formation-time information set are
used for both models. This restriction makes the model contrast interpretable:
differences between the two forecasting methods cannot be attributed simply
to one model receiving additional firm characteristics. However, it also
means that any apparent nonlinear advantage must be interpreted relative
to a sparse baseline; important known predictors and stronger simple
characteristic benchmarks are reserved for explicit robustness analyses.

\subsection{Volatility, market states, and conditional predictability}
\label{subsec:literature-volatility}

Market volatility summarizes more than a single structural mechanism.
\cite{BekaertHoerova2014} show that the VIX contains information connected
to expected variance and variance risk premia. The relation between
variance risk premia and subsequent aggregate stock-market returns is
studied by \cite{BollerslevTauchenZhou2009}. These contributions help
motivate a VIX-based conditioning variable, but their aggregate-return
outcomes are not identical to individual-stock return rankings.
A VIX state is also observational: it co-moves with risk perceptions,
funding conditions, liquidity and the business cycle, and it is not an
exogenously assigned intervention.

Studies of volatility-managed factor portfolios illustrate that an investor's
exposure to time-varying volatility may matter even if unconditional factor
relationships remain the same \cite{MoreiraMuir2017}.
Momentum provides a particularly relevant example because it belongs to
the current study's frozen characteristic set:
\cite{DanielMoskowitz2016} examine momentum crashes associated with
panic-like market conditions, and \cite{BarrosoSantaClara2015}
investigate changing momentum risk. A high-volatility forecast-performance
difference could therefore reflect an interaction between VIX states and
momentum exposure rather than a general property of nonlinear machine
learning.

The study uses two historically defined, but conceptually distinct,
classifications. The original binary LOW/HIGH scheme compares formation-month
VIX with its expanding historical median. The additional LOW/MIDDLE/HIGH
evaluation uses expanding historical VIX terciles. Both are calculated using
information available at portfolio formation, and neither involves fitting
separate models in each regime. Their thresholds are reproducible empirical
conventions rather than established structural break points. Consequently,
alternative VIX thresholds and realized-volatility measures are necessary
robustness checks before drawing claims about the market-state nature of
predictability.

\subsection{Conditional nonlinear predictability and closest competing studies}
\label{subsec:literature-closest}

A major originality challenge is that volatility-dependent machine-learning
returns and implementation constraints have already been studied directly.
\cite{AvramovChengMetzker2023} find that investment returns from deep-learning
signals are concentrated in difficult-to-arbitrage stocks and deteriorate
when microcapitalization stocks, distressed firms, or high-volatility
episodes are excluded; turnover and trading costs further reduce
profitability. This substantially overlaps the motivation of the present
analysis. It would therefore be incorrect to claim that linking machine
learning, market volatility, microcaps and transaction costs is itself
unprecedented.

The current empirical design differs in several more specific dimensions.
It evaluates historically frozen forecasts from \emph{two} model classes on
the same three-characteristic information set; partitions the same monthly
forecast record using historically observable binary and three-state VIX
classifications; distinguishes stock-level forecast loss from monthly
cross-sectional rank information; tests \emph{direct} differences between
regime-conditional means with heteroskedasticity- and
autocorrelation-consistent inference; and explicitly compares gross and
cost-adjusted equal- and value-weighted decile-spread portfolios.
A recent non-refereed study by \cite{FangSlepaczuk2026} likewise
combines regime-aware volatility forecasting, XGBoost return prediction,
walk-forward validation and transaction-cost discussion. Its Chinese
high-frequency index setting differs from the present monthly
individual-stock cross-section, but the overlapping research question
further limits broad claims of priority.

Those design features can improve interpretability. Nevertheless, assembling
existing methods into one pipeline is not by itself proof of a sufficiently
distinct journal contribution. That determination requires a full-text
comparison with competing studies and further empirical evidence that resolves
a substantive question not already answered.

The literature also warns against attributing all state dependence to
the model architecture. \cite{KellyPruittSu2019} provide a risk-exposure
interpretation of time-varying characteristic relations, while
\cite{DanielMoskowitz2016} highlight the potentially nonlinear downside
risk embedded in momentum. A model may rank stocks differently because
its characteristic responses are nonlinear, because the composition of the
eligible cross-section changes, or because macroeconomic exposures vary by
state. These explanations should be separated as far as observable data
permit, but the current design cannot establish a causal channel.

\subsection{Statistical evaluation, economic value, and investment frictions}
\label{subsec:literature-evaluation}

The evaluation of forecasts requires matching the statistic to the
research question. Mean squared forecasting error, out-of-sample
$R^2$, a monthly Spearman rank correlation, and a realized decile spread
capture different properties of a prediction. Evidence of modest
out-of-sample improvement for aggregate equity-premium forecasts
\cite{CampbellThompson2008,WelchGoyal2008,RapachStraussZhou2010}
provides useful caution about benchmarks and stability, but should not
be presented as direct evidence about the cross-section of individual
securities.

Comparing two forecast models on a common set of realized outcomes
also creates dependence. \cite{DieboldMariano1995} and
\cite{GiacominiWhite2006} provide influential approaches to comparing
predictive accuracy and conditional predictive ability, respectively,
under their stated sampling and estimation assumptions. Those results do
not license treating hundreds of thousands of stock--month residuals
as independent test observations. In this paper the formation month is
the common decision and conditioning unit. Additional model-comparison
tests will therefore form aligned monthly loss or rank-difference series
before applying an appropriate dependent-observation inference procedure.

For the existing monthly mean-return, rank-IC and regime-difference analyses,
the reported framework uses the heteroskedasticity- and
autocorrelation-consistent covariance estimator associated with
\cite{NeweyWest1987}. Direct LOW/MIDDLE/HIGH pairwise contrasts
are evaluated with the family-wise step-down adjustment of
\cite{Holm1979}; the false-discovery-rate procedure of
\cite{BenjaminiHochberg1995} is a separately identified secondary
diagnostic. This distinction is important: a positive within-LOW estimate,
together with an insignificant within-HIGH estimate, does not constitute
a significant LOW-minus-HIGH difference. The direct contrast, its uncertainty,
and any specified multiple-testing correction must be reported.

Economic evaluation raises further independent challenges.
More recently, \cite{JensenKellyMalamudPedersen2026} develop an
implementable efficient frontier in which machine-learning portfolio
weights are learned with trading costs explicitly incorporated into
the economic objective. Their research shows why simple forecast-then-
trade policies can fail to deliver net performance when strategies
rely on fleeting small-stock characteristics. This is a different
optimization objective from the current frozen decile portfolios,
but it establishes that cost-aware investment use of machine learning
is already a substantial published research area.


\cite{NovyMarxVelikov2016} document the economic consequences of trading
costs across anomalies, while \cite{DeMiguelEtAl2020} analyze the influence
of transaction costs when exploiting multiple stock characteristics.
\cite{HouXueZhang2020} emphasize the sensitivity of anomaly evidence
to replication design, microcaps and statistical testing.
\cite{HarveyLiuZhu2016} caution that multiple testing complicates
interpretation of large collections of predictive findings, and
\cite{McLeanPontiff2016} study attenuation of anomaly returns across
subsequent and post-publication periods.

The present portfolio analysis sorts next-month predicted returns into
deciles and evaluates $D10-D1$ with equal and lagged-market-equity value
weights. The cost-adjusted three-state results are based on portfolio
turnover from return-drifted holdings and a stipulated one-way cost of
50 basis points per traded dollar for both long and short legs.
This is a reproducible \emph{scenario}, not an estimate of actual
execution, borrowing, financing or market-impact costs. In particular,
positive modeled net returns are insufficient to establish feasible
capacity or implementable profitability.

\subsection{Potential research gap and conditional contribution}
\label{subsec:literature-gap}

Taken together, the preceding studies leave a possible, but not yet
demonstrated, narrower contribution: a unified assessment of how
cross-sectional \emph{ranking quality}, \emph{relative model accuracy},
and \emph{net portfolio value} behave across historically observable
VIX states when both models use an identical, transparent baseline
information set and direct regime comparisons are subjected to a common
inference protocol. The distinction between descriptive within-state
performance and formal between-state heterogeneity is central.

This framing is deliberately conditional on the results already obtained.
The production baseline exhibits little overall next-month stock-level
forecast-error improvement relative to its historical-mean benchmark.
The original binary classification shows positive rank information and
a positive modeled net equal-weighted spread for XGBoost in LOW months,
but direct LOW/HIGH differences remain imprecise.
In the three-state evaluation, XGBoost's equal-weighted cost-adjusted
mean spread is approximately $0.614\%$ per month for LOW and
$0.960\%$ for MIDDLE, compared with $-0.515\%$ for HIGH.
However, the three-state net pairwise differences do not survive the
specified Holm adjustment, and the corresponding value-weighted net
spreads are negative. Thus, the existing results do not establish a
statistically significant ordering of regimes or economically
implementable profits. Distinct claims about stronger benchmarks,
threshold robustness, microcap attribution and interpretation await
additional analysis.

A negative outcome in those subsequent exercises need not be concealed.
If state-conditional differences disappear under better benchmarks,
historically justified volatility thresholds, investability restrictions,
or robust sample splits, the result would delimit the reliability
of simple characteristic-based machine-learning strategies. Whether such
an outcome contributes beyond existing evidence remains a question
for the final competitor audit and faculty review.

\subsection{Research questions and transparent hypothesis framework}
\label{subsec:literature-hypotheses}

Let $S_t\in\{L,M,H\}$ denote the three-state historical VIX label of
formation month $t$. For forecasting model $m$ define
\[
IC_t^{m}=\mathrm{Spearman}_{i}
 \left(\widehat R_{i,t+1}^{m},R_{i,t+1}\right),
\]
and let $G_{w,t}^{m}$ and $N_{w,t}^{m}$ respectively denote gross and
50-basis-point modeled net returns on its $D10-D1$ strategy with
weighting convention $w\in\{EW,VW\}$. These monthly outcomes are
evaluated conditional on the formation-month state, not by
training models separately within states.

\paragraph{H1: Regime heterogeneity.}
For a declared outcome $Y_t^m$ (starting with monthly rank IC and
separately considering gross and net portfolio returns), the null is
\[
H_{0,1}:\quad
 E[Y_t^m\mid S_t=L]=E[Y_t^m\mid S_t=M]
 =E[Y_t^m\mid S_t=H].
\]
The two-sided alternative is that at least one conditional mean differs.
Joint time-series tests and multiplicity-controlled pairwise contrasts,
rather than isolated within-state significance, address this question.

\paragraph{H2: Relative nonlinear predictive advantage.}
Let
\[
\Delta_r^{IC}
 = E[IC_t^{XGB}-IC_t^{EN}\mid S_t=r].
\]
The state-heterogeneity null is
$H_{0,2}:\Delta_L^{IC}=\Delta_M^{IC}=\Delta_H^{IC}$,
with a corresponding null for aligned monthly forecast-loss
differences, where the loss sign is specified before estimation.
The relevant object is a \emph{paired} model difference, not
a comparison of two separately computed significance stars.

\paragraph{H3: Conditional economic value.}
The observed three-state net equal-weighted LOW XGBoost mean motivates
an examination of whether the conditional expected modeled net
$D10-D1$ return is positive and remains so under investability filters.
A positive sample mean, a statistically nonzero mean, factor-adjusted
performance and executable profit are not equivalent conclusions.
Tests already run on the current sample are retrospective evidence.

\paragraph{H4: Intermediate-volatility maximum (exploratory).}
A peak requires both
$E[Y_t\mid S_t=M]>E[Y_t\mid S_t=L]$ and
$E[Y_t\mid S_t=M]>E[Y_t\mid S_t=H]$
for a predeclared outcome. This possibility arose after the
three-state results were inspected. It is therefore exploratory, not
an ex ante confirmatory claim, and requires an appropriately joint
or simultaneous test plus additional independent evidence.

\paragraph{H5: Investability and stability.}
Future tests will evaluate whether any positive conditional net
spread remains after prespecified microcapitalization, liquidity,
borrow-cost and other implementability restrictions.
This is a future analysis plan, not established evidence from the
present cost scenario.

The full hypothesis registry, including estimands, dependence assumptions,
test families, decision rules, exploratory status and mappings to future
research milestones, is maintained in the project's Milestone 11
documentation. All hypotheses are documented after the M1--M10
results were observed. This section must not be interpreted as
retrospective preregistration.